\documentclass[../../../main.tex]{subfiles}
\begin{document}

\subsection{Angular Momentum of Single Particle}
From the definition $\mathbf{L}= \mathbf{r} \times \mathbf{p}$, we can obtain the quantum angular momentum
\begin{equation*}
    \hat{L}_i = \epsilon_{ijk}r_jp_k
\end{equation*}
The commutation relation for angular are obtained from position $\mathbf{r}$ and momentum $\mathbf{p}$ relation $[r_i   , p_j]=i \hbar \delta_{ij}$
\begin{equation*}
    [\hat{L}_{i}, \hat{L}_{j}] = i \hbar \epsilon_{ijk} \hat{L}_{k}
\end{equation*}
where $\varepsilon_{jkl}$ are the components of the totally antisymmetric Levi-Civita tensor
\begin{equation*}
    \epsilon_{jkl} = \begin{cases}
        +1 & \text{if } (jkl) \text{ is an even permutation of } (123) \\
        -1 & \text{if } (jkl) \text{ is an odd permutation of } (123)  \\
        0  & \text{if any two indices are equal}
    \end{cases}
\end{equation*}
Even permutations of $(123)$, beside the identity $(123)$ itself, are
\begin{align*}
    (231) & =  \mathcal{T}_{13} \mathcal{T}_{12}(123) \\
    (312) & =  \mathcal{T}_{12} \mathcal{T}_{13}(123)
\end{align*}
while the odd permutations are
\begin{align*}
    (213) & =  \mathcal{T}_{12}(123) \\
    (132) & =  \mathcal{T}_{23}(123) \\
    (321) & =  \mathcal{T}_{13}(123)
\end{align*}

Another one for the position case
\begin{equation*}
    [r_{i}, \hat{L}_{j}] = i \hbar \epsilon_{ijk} r_{k}
\end{equation*}
and momentum
\begin{equation*}
    [p_{i}, \hat{L}_{j}] = i \hbar \epsilon_{ijk} p_{k}
\end{equation*}

It follows from the commutator relation and the corresponding Heisenberg uncertainty relation that the components $\hat{L}_j$ of the angular momentum operator $L$ cannot be sharply measured simultaneously.
However, for the square of the angular momentum $L^2$
\begin{equation*}
    [L^2,\hat{L}_j] = 0
\end{equation*}
we find that the square $L^2$, and one component $Lj$ of the angular momentum $L$, can be measured simultaneously with arbitrary precision.

The simultaneous eigenstates of $\hat{L}^2$
\begin{equation*}
    \hat{L}^2 \ket{l,m} = l \hbar ^2 (l+1) \ket{l,m}
\end{equation*}
and $\hat{L}_3$
\begin{equation*}
    \hat{L }_3 \ket{m}= m \hbar \ket{m}
\end{equation*}
is denoted by
\begin{equation*}
    \left| l, m \right\rangle; \text{ such that } 2l = 0, 1, 2, \ldots, \quad -l \leq m \leq l
\end{equation*}
The eigenvalue equation of $\hat{L}^2$ shows that $l$ determines the magnitude of the angular momentum vector, specifically $|\mathbf{L}|=\hbar \sqrt{l(l+1)}$.
In other hand, the eigenvalue equation of $\hat{L }_3$ measures the component of angular momentum along a chosen quantization axis, that is usually chosen to be $z$ axis.

The annihilation and  creation operators of angular momentum are expressed as
\begin{equation*}
    \hat{L}_{\pm} = \hat{L}_1 \pm i \hat{L}_2
\end{equation*}
or equivalently
\begin{equation*}
    \hat{L}_1 = \frac{1}{2}(\hat{L}_{+} + \hat{L}_{-})\qquad
    \hat{L}_2 = \frac{1}{2i}(\hat{L}_{+} - \hat{L}_{-})
\end{equation*}
These satisfy the commutator relation
\begin{equation*}
    [\hat{L}_3,\hat{L}_{\pm}] = \pm \hbar  \hat{L}_{\pm}, \qquad
    [\hat{L}_+,\hat{L}_-]=2\hat{L}_3
\end{equation*}
Their action on eigenstates $\ket{l,m}$ are
\begin{equation*}
    \hat{L}_{\pm}   |l, m\rangle = \hbar \sqrt{l(l+1) - m(m \pm 1)}   |l, m \pm 1\rangle
\end{equation*}

For $l=1/2$, each orbital momentum can be described by $2 \times 2$ matrix
\begin{equation*}
    \hat{L}_1 = \frac{\hbar}{2}\begin{bmatrix} 0 & 1 \\ 1 & 0 \end{bmatrix},\quad
    \hat{L}_2 = \frac{\hbar}{2}\begin{bmatrix} 0 & -i \\ i & 0 \end{bmatrix},\quad
    \hat{L}_3 = \frac{\hbar}{2}\begin{bmatrix} 1 & 0 \\ 0 & -1 \end{bmatrix}
\end{equation*}
as for the annihilation and creation
\begin{equation*}
    \hat{L}_{+} = \hbar \begin{bmatrix} 0 & 1 \\ 0 & 0 \end{bmatrix}\qquad
    \hat{L}_{-} = \hbar \begin{bmatrix} 0 & 0 \\ 1 & 0 \end{bmatrix}
\end{equation*}
For $l=1$, they are $3 \times 3$
\begin{equation*}
    \hat{L}_1 = \frac{\hbar}{\sqrt{2}} \begin{bmatrix} 0 & 1 & 0 \\ 1 & 0 & 1 \\ 0 & 1 & 0 \end{bmatrix},\;
    \hat{L}_2 = \frac{\hbar}{\sqrt{2}i} \begin{bmatrix} 0 & 1 & 0 \\ -1 & 0 & 1 \\ 0 & -1 & 0 \end{bmatrix},\;
    \hat{L}_3 = \hbar \begin{bmatrix} 1 & 0 & 0 \\ 0 & 0 & 0 \\ 0 & 0 & -1 \end{bmatrix}
\end{equation*}
as for the annihilation and creation
\begin{equation*}
    \hat{L}_{+} = \hbar \begin{bmatrix} 0 & \sqrt{2} & 0 \\ 0 & 0 & \sqrt{2} \\ 0 & 0 & 0 \end{bmatrix}\qquad
    \hat{L}_{-} = \hbar \begin{bmatrix} 0 & 0 & 0 \\ \sqrt{2} & 0 & 0 \\ 0 & \sqrt{2} & 0 \end{bmatrix}
\end{equation*}
Note that the matrices are written in the $\ket{l,m}$ basis ordered by descending $m$, that is
\begin{equation*}
    \ket{\psi} =
    \begin{bmatrix}
        \psi_{m=l} \\ \psi_{m=l-1} \\ \vdots \\ \psi_{m=-l}
    \end{bmatrix}
\end{equation*}

\subsubsection*{Derivation: angular momentum components commutator relation.}
Before we begin, consider the identities of commutator
\begin{align*}
    [\mathrm{A},\mathrm{B}\mathrm{C}] & =\mathrm{B}[\mathrm{A},\mathrm{C}]+[\mathrm{A},\mathrm{B}]\mathrm{C} \\
    [\mathrm{A}\mathrm{B},\mathrm{C}] & =\mathrm{A}[\mathrm{B},\mathrm{C}]+[\mathrm{A},\mathrm{C}]\mathrm{B} \\
    [\mathrm{A}\mathrm{B},\mathrm{C}\mathrm{D}]
    [\mathrm{A}\mathrm{B},\mathrm{C}\mathrm{D}]
                                      & = \mathrm{A}[\mathrm{B},\mathrm{C}]\mathrm{D}
    + \mathrm{A}\mathrm{C}[\mathrm{B},\mathrm{D}]
    + [\mathrm{A},\mathrm{C}]\mathrm{D}\mathrm{B}
    + \mathrm{C}[\mathrm{A},\mathrm{D}]\mathrm{B}
\end{align*}

The $i$-th component of cross product is written as
\begin{equation*}
    (A \times B)_i = \epsilon_{ijk} A_j B_k
\end{equation*}
Hence from the definition of angular momentum $\mathbf{L }= \mathbf{r }\times \mathbf{p}$, we can write
\begin{equation*}
    \hat{L}_i =\epsilon_{ijk}r_jp_k
\end{equation*}
Using this we can derive the commutator relation
\begin{align*}
    [\hat{L}_i, \hat{L}_j]
     & = \epsilon_{ilm} \epsilon_{jnp} [ r_lp_m,r_np_p]                                                                                                                             \\
     & =  \epsilon_{ilm} \epsilon_{jnp}(  r_l[p_m, r_n]p_p + r_l r_n[p_m, p_p] + [r_l, r_n] p_m p_p                                                                                 \\
     & \qquad  + r_n [r_l, p_p] p_m)                                                                                                                                                \\
     & =  \epsilon_{ilm} \epsilon_{jnp}\left( - i \hbar  r_l \delta_{mn} p_p + i \hbar  r_n \delta_{lp} p_m \right)                                                                 \\
     & = i \hbar \left( -\epsilon_{ilm }\epsilon_{jmp }r_lp_p + \epsilon_{ilm }\epsilon_{jnl }  r_np_m\right)                                                                       \\
     & = i \hbar \left( \epsilon_{lmi }\epsilon_{ljn}r_n  p_m - \epsilon_{mil }\epsilon_{mpj }r_l p_p\right)                                                                        \\
     & = i \hbar \left[ \left( \delta_{mj }\delta_{in }-\delta_{mn }\delta_{ji } \right)r_np_m - \left( \delta_{ip }\delta_{lj } - \delta_{ij }\delta_{pl }\right)r_lp_p    \right] \\
     & = i \hbar \left[ \left( r_i p_j -\delta_{ij }r_m p_m\right) -\left( r_jp_i-\delta_{ij }r_l p_p \right)   \right]                                                             \\
     & = i \hbar \left[ r_i p_j  -r_jp_i \right]                                                                                                                                    \\
    [\hat{L}_i, \hat{L}_j]
     & = i \hbar \epsilon_{ijk}\hat{L}_k
\end{align*}
Next for the commutator with posistion
\begin{align*}
    [r_i,\hat{L}_j]
     & =  \epsilon_{jnp}[r_i,r_np_p]                                       \\
     & = \epsilon_{jnp} \left( r_n[r_i,p_p ]+[r_i,r_n ]p_p \right)         \\
     & =  i \hbar \epsilon_{jnp }r_n \delta_{ip}=i \hbar \epsilon_{jni}r_n \\
    [r_i,\hat{L}_j]
     & = i \hbar \epsilon_{jki}r_k = i \hbar \epsilon_{ijk}r_k
\end{align*}
Finally for the linear momentum
\begin{align*}
    [p_i,\hat{L}_j]
     & =  \epsilon_{jnp}[p_i,r_np_p]                                           \\
     & = \epsilon_{jnp} \left( r_n[p_i,p_p ]+[p_i,r_n ]p_p \right)             \\
     & =  -i \hbar \epsilon_{jnp }p_p \delta_{in} = i \hbar  \epsilon_{ijp}p_p \\
    [p_i,\hat{L}_j]
     & = i \hbar \epsilon_{ijk}p_k
\end{align*}

\subsubsection*{Derivation: angular momentum square commutator relation.}
Derived as follows
\begin{align*}
    [L^2, \hat{L}_j]
     & = [\hat{L}_i \hat{L}_i, \hat{L}_j] =  \hat{L}_i[\hat{L}_i, \hat{L}_j] + [\hat{L}_i, \hat{L}_j]\hat{L}_i \\
    [L^2, \hat{L}_j]
     & = i \hbar \epsilon_{ijk} (\hat{L}_i \hat{L}_k + \hat{L}_k \hat{L}_i)
\end{align*}
If we define $S_{ik}=\hat{L}_i \hat{L}_k + \hat{L}_k \hat{L}_i$, which symmetric under $i \leftrightarrow k$, notice that by performing $i \leftrightarrow k$ transposition
\begin{align*}
    \epsilon_{ijk}S_{ik} & = \epsilon_{kji}S_{ki}  \\
    \epsilon_{ijk}S_{ik} & = -\epsilon_{ijk}S_{ik} \\
    \epsilon_{ijk}S_{ik} & = 0
\end{align*}
Hence the commutator evaluate into zero.

\subsubsection{Derivation: annihilation and creation operator commutator.}
For the commutator of $\hat{L}_3$ with annihilation and creation
\begin{align*}
    [\hat{L}_3, \hat{L}_\pm]
     & =  [\hat{L}_3, \hat{L}_1 \pm i\hat{L}_2] = [\hat{L}_3, \hat{L}_1] \pm i[\hat{L}_3, \hat{L}_2] \\
     & =  i\hbar \hat{L}_2 \pm i(-i\hbar \hat{L}_1)=  i\hbar \hat{L}_2 \pm \hbar \hat{L}_1           \\
     & = \begin{cases}
             \hbar(\hat{L}_1 + i\hat{L}_2) = \hbar \hat{L}_+,  & \text{for } + \\
             -\hbar(\hat{L}_1 - i\hat{L}_2) = -\hbar \hat{L}_- & \text{for } - \\
         \end{cases}                           \\
     & =  \pm \hbar \hat{L}_\pm
\end{align*}
While for the commutator of annihilation and creation with each other
\begin{align*}
    [\hat{L}_{+}, \hat{L}_{-}]
     & =  [\hat{L}_1 + i\hat{L}_2, \hat{L}_1 - i\hat{L}_2]                                                           \\
     & =  [\hat{L}_1, \hat{L}_1] - i[\hat{L}_1, \hat{L}_2] + i[\hat{L}_2, \hat{L}_1] + (i)(-i)[\hat{L}_2, \hat{L}_2] \\
     & =  -i[\hat{L}_1, \hat{L}_2] + i[\hat{L}_2, \hat{L}_1]=-2i[\hat{L}_1, \hat{L}_2]                               \\
    [\hat{L}_{+}, \hat{L}_{-}]
     & =  2\hbar \hat{L}_3                                                                                           \\
\end{align*}

\subsubsection{Derivation: annihilation and creation operator action.}
With $\ket{m}$ as the eigenstate of $\hat{L}_3$, we define shift operator $\hat{L}_S$ which produce new state $\ket{\lambda}$ for this eigenstate
\begin{equation*}
    \ket{\lambda}=\hat{L}_S \ket{m}
\end{equation*}
The action of $\hat{L}_3$ is thus
\begin{align*}
    \hat{L}_3 |\lambda\rangle       & =  (m+\lambda) \hbar \ket{\lambda}           \\
    \hat{L}_3 (\hat{L}_S |m\rangle) & =  (m + \lambda) \hbar (\hat{L}_S |m\rangle)
\end{align*}
Multiplying the eigenvalue equation of $\hat{L}_3$ with $\hat{L}_S$
\begin{equation*}
    \hat{L}_S\hat{L}_3 \ket{m}=\hat{L}_Sm \hbar \ket{m}
\end{equation*}
then substracting it to the $\ket{\lambda}$ equation
\begin{align*}
    (\hat{L}_3\hat{L}_S - \hat{L}_S\hat{L}_3) \ket{m} & = \lambda \hat{L}_S \ket{m} \\
    [\hat{L}_3,\hat{L}_S]                             & = \lambda \hat{L}_S
\end{align*}
For $\hat{L}_S$ to satisfy this relation, we can start by assuming that it consists as linear combination
\begin{equation*}
    \hat{L}_S = a\hat{L}_1 + b\hat{L}_2
\end{equation*}
We have then the commutator relation
\begin{equation*}
    [\hat{L}_3, \hat{L}_S] = a[\hat{L}_3, \hat{L}_1] + b[\hat{L}_3, \hat{L}_2] = a(i\hbar \hat{L}_2) + b(-i\hbar \hat{L}_1) = -i\hbar b \hat{L}_1 + i\hbar a \hat{L}_2
\end{equation*}
and the shift condition
\begin{equation*}
    [\hat{L}_3, \hat{L}_S] = \lambda \hat{L}_S = \lambda (a \hat{L}_1 + b \hat{L}_2)
\end{equation*}
By equating them
\begin{align*}
    \lambda (a \hat{L}_1 + b \hat{L}_2) & =  -i\hbar b \hat{L}_1 + i\hbar a \hat{L}_2 \\
    (\lambda a + i \hbar b)\hat{L}_1    & =  (i \hbar a - \lambda b )\hat{L}_2
\end{align*}
We obtained linear system
\begin{equation*}
    \lambda a + i \hbar b =0 \qquad
    i \hbar a - \lambda b   =0
\end{equation*}
or in matrix representation
\begin{equation*}
    \begin{bmatrix}
        \lambda & i \hbar  \\
        i \hbar & -\lambda
    \end{bmatrix}
    \begin{bmatrix}
        a \\
        b
    \end{bmatrix}
    =
    \begin{bmatrix}
        0 \\
        0
    \end{bmatrix}
\end{equation*}
As is the case with eigenvalue problem, the determinant of coefficient matrix must equal to zero for the non-trivial solution to exist.
Hence
\begin{equation*}
    -\lambda^2+\hbar ^2=0\implies \lambda=\pm \hbar
\end{equation*}
The shift operator can thus be identified with the operators $L^+$ and $L^-$, creating states where the eigenvalue $m$ of an eigenstate $|m\rangle$ of $\hat{L}_3$ is shifted by $\pm \hbar $
\begin{equation*}
    \hat{L}_3 (\hat{L}_S |m\rangle) = (m \pm 1)\hbar  (\hat{L}_S |m\rangle)
\end{equation*}

We consider next a simultaneous eigenstate $|\Psi\rangle$ of $L^2$ and $\hat{L}_3$ with eigenvalues $C^2$ and $m$, respectively
\begin{equation*}
    L^2|\Psi\rangle = C^2|\Psi\rangle     \qquad
    \hat{L}_3|\Psi\rangle = m \hbar |\Psi\rangle
\end{equation*}
which, as we have seen, is always possible.
The action of  $L^\pm$ on this state can be seen by considering the states $|\Phi^\pm\rangle = L^\pm|\Psi\rangle$.
We already know that the eigenvalue of $\hat{L}_3$ gets shifted
\begin{equation*}
    \hat{L}_3|\Phi^\pm\rangle = (m \pm 1) \hbar |\Phi^\pm\rangle
\end{equation*}
More generally, we have, by iteration,
\begin{equation*}
    \hat{L}_3 \left( (L^{\pm})^n  \ket{\Psi} \right) = (m \pm n) \hbar  \left( (L^{\pm})^n \ket{\Psi} \right)
\end{equation*}
However
\begin{equation*}
    L^2 \left( (L^{\pm})^n | \Psi \right) = (L^{\pm})^n L^2 \ket{\Psi} = C^2 \left( (L^{\pm})^n | \Psi \right)
\end{equation*}
The  expectation values in this state are
\begin{equation*}
    \braket{\Phi | \hat{L}_3|\Phi} = (m+n) \hbar \qquad
    \braket{\Phi | L^2|\Phi} = C^2
\end{equation*}
By definition $L^2=\hat{L}_1^2+\hat{L}_2^2+\hat{L}_3^2$, the square of the third component of the angular momentum cannot exceed the square of the angular momentum
\begin{equation*}
    (m + n)^2 \hbar ^2 \leq C^2
\end{equation*}
Thus, \( m \) must be bounded from above and below
\begin{equation*}
    m_{\text{max}} \leq m \leq m_{\text{min}}
\end{equation*}
such that
\begin{equation*}
    \hat{L}_3 | m_{\text{max}} \rangle = m_{\text{max}} \hbar  | m_{\text{max}} \rangle \qquad
    L^+ | m_{\text{max}} \rangle = 0
\end{equation*}
and
\begin{equation*}
    \hat{L}_3 | m_{\text{min}} \rangle = m_{\text{min}} \hbar  | m_{\text{min}} \rangle \qquad
    L^- | m_{\text{min}} \rangle = 0
\end{equation*}
Writing $L^2$ as
\begin{align*}
    L^2 & = \frac{1}{4} \left( \hat{L}_+^2 + \hat{L}_+ \hat{L}_- + \hat{L}_- \hat{L}_+ + \hat{L}_-^2 \right)                     \\
        & \qquad - \frac{1}{4} \left( \hat{L}_+^2 - \hat{L}_+ \hat{L}_- - \hat{L}_- \hat{L}_+ + \hat{L}_-^2 \right) +\hat{L}_3^2 \\
    L^2
        & = \hat{L}_3^2 +\hbar  \hat{L}_3 + L^- L^+ =  \hat{L}_3^2 - \hat{L}_3 + L^+ L^-
\end{align*}
using $    [\hat{L}_+,\hat{L}_-]=2 \hbar \hat{L}_3$.
Choosing $| m_{\text{max}} \rangle$ and $| m_{\text{min}} \rangle$ as simultaneous eigenstates of $L^2$, we find
\begin{align*}
    L^2 | m_{\text{max}} \rangle & =   \left( L^- L^+ +  \hat{L}_3^2 +\hbar  \hat{L}_3 \right) | m_{\text{max}} \rangle \\
    C^2 | m_{\text{max}} \rangle
                                 & =  m_{\text{max}} \hbar ^2(m_{\text{max}} + 1) | m_{\text{max}} \rangle
\end{align*}
and
\begin{align*}
    L^2 |m_\text{min}\rangle & = \left( L^+ L^- + L^2_3 -\hbar  \hat{L}_3 \right) |m_\text{min}\rangle \\
    C^2 |m_\text{min}\rangle
                             & =  m_\text{min} \hbar ^2 (m_\text{min} - 1) |m_\text{min}\rangle
\end{align*}
Equating both, we have
\begin{align*}
    m_{\text{max}} (m_{\text{max}} + 1) & =   m_\text{min} (m_\text{min} - 1) \\
    m^2_{\text{max}} + m_{\text{max}}   & =  m^2_{\text{min}} -_{\text{min}}
\end{align*}
and
\begin{align*}
    m^2_{\text{max}} - m^2_{\text{min}} + m_{\text{max}} + m_{\text{min}}                                  & = 0  \\
    (m_{\text{max}} - m_{\text{min}})(m_{\text{max}} + m_{\text{min}}) + (m_{\text{max}} + m_{\text{min}}) & = 0  \\
    (m_{\text{max}} + m_{\text{min}})\left[(m_{\text{max}} - m_{\text{min}}) + 1\right]                    & =  0
\end{align*}
Thus two possibilities exits: $m_\text{min} = -m_\text{max}$, which is the only admissible solution, and $m_\text{max} < m_\text{min}$,which contradicts the definition of maximum and minimum.

Denoting  $m_{\text{max}}$ by the letter $l$, and the corresponding simultaneous eigenstate $\ket{l,l}$ of  $L^2$ and $\hat{L}_3$  satisfies
\begin{equation*}
    L^+ |l, l\rangle = 0\quad
    \hat{L}_3 |l, l\rangle = l \hbar  |l, l\rangle\quad
    L^2 |l, l\rangle  = l h^2(l + 1) |l, l\rangle\quad
\end{equation*}
where the notation $\lvert l, l \rangle$ simply indicates the eigenvalues of $L^2$ (first slot) and $\hat{L}_3$ (second slot).
Moreover, we can conclude from the action of $\hat{L}_+$ to the state $\ket{l,l}$, repeated $n$ application of $\hat{L}_-$ gives $m=l-n$.
The lowest value occure  at
\begin{equation*}
    -l= l-n \implies 2l=n
\end{equation*}
This mean  $l$ can only take integer or half-odd integer values.

So far, the action of $\hat{L}_\pm$ can be written as
\begin{equation*}
    L^{\pm} |l,m\rangle = N^{\pm}_{l,m} |l,m \pm 1\rangle
\end{equation*}
assuming that both states $|l,m\rangle$ and $|l,m \pm 1\rangle$ are normalized.
All that left us to determine the normalization constant
\begin{align*}
    |N^{\pm}_{l,m}|^2
     & = \langle l,m|L^{\mp}L^{\pm}|l,m\rangle                              \\
     & =  \langle l,m| (L^2 - \hat{L}_3^2 \mp \hbar  \hat{L}_3) |l,m\rangle \\
     & =  l \hbar^2 (l + 1) - m^2 \hbar ^2 \mp m \hbar                      \\
    N^{\pm}_{l,m}
     & = \hbar \sqrt{l(l + 1) - m(m \pm 1)}
\end{align*}

\subsection{Spin }
There is a fundamental distinction between particles with half-odd integer spin, called Fermions, which obey the Pauli exclusion principle, and those with integer spin, called Bosons, which do not obey the Pauli exclusion principle.
For special case $l=1/2\equiv s$, we introduce operator $\mathbf{S}$, which satisfies
\begin{equation*}
    [\hat{S}_{i}, \hat{S}_{j}] = i \epsilon_{ijk} \hat{S}_{k}
\end{equation*}
The weight of a state $|\chi\rangle$ measures the eigenvalue of the spin operator of the state of the quantum particle
\begin{equation*}
    S_3 |\chi\rangle = m_S |\chi\rangle = \pm \frac{1}{2} |\chi\rangle \quad
    S^2 |\chi\rangle = s(s + 1) |\chi\rangle = \frac{3}{4} |\chi\rangle
\end{equation*}
with
\begin{equation*}
    -s \leq m_S \leq s, \quad  m_S = \pm \frac{1}{2}
\end{equation*}
where $s$ denote the spin of quantum particle.

The general element of the space spawned by this basis is called a spinor and can be formally written as
\begin{equation*}
    \ket{\boldsymbol{\chi}}= \begin{bmatrix}
        \ket{\chi_+} \\
        \ket{\chi_-}
    \end{bmatrix}
\end{equation*}
In this coordinate system, we have two possible state: one for $m_s=1/2$ and other for $m_s=-1/2$
\begin{equation*}
    \big| \frac{1 }{2 },\frac{1 }{2 }\big\rangle = \ket{\uparrow} = \begin{bmatrix} 1\\0 \end{bmatrix}\qquad
    \big| \frac{1 }{2 },-\frac{1 }{2 }\big\rangle = \ket{\downarrow} = \begin{bmatrix} 0\\1 \end{bmatrix}
\end{equation*}

The components of spin $S_i$ can be written in terms of Pauli spin matrix
\begin{equation*}
    \mathbf{S} = \frac{1 }{2 } \boldsymbol{\sigma}
\end{equation*}
with
\begin{equation*}
    \sigma_1 =
    \begin{pmatrix}
        0 & 1 \\
        1 & 0
    \end{pmatrix}
    \quad
    \sigma_2 =
    \begin{pmatrix}
        0 & -i \\
        i & 0
    \end{pmatrix}
    \quad
    \sigma_3    =
    \begin{pmatrix}
        1 & 0  \\
        0 & -1
    \end{pmatrix}
\end{equation*}
such that
\begin{equation*}
    S_1 = \frac{1}{2}
    \begin{pmatrix}
        0 & 1 \\
        1 & 0
    \end{pmatrix}
    \quad
    S_2 = \frac{1}{2}
    \begin{pmatrix}
        0 & -i \\
        i & 0
    \end{pmatrix}
    \quad
    S_3 = \frac{1}{2}
    \begin{pmatrix}
        1 & 0  \\
        0 & -1
    \end{pmatrix}
\end{equation*}
The Pauli matrices satisfy
\begin{gather*}
    \sigma_i^2=\mathbf{1}  \\
    [\sigma_i, \sigma_j] = 2i \epsilon_{ijk} \sigma_k\qquad    \left\{ \sigma_i,\sigma_j \right\} = i\delta_{ij }\mathbf{1}\\
    (\mathbf{a} \cdot \boldsymbol{\sigma})(\mathbf{b} \cdot \boldsymbol{\sigma}) = \mathbf{a} \cdot \mathbf{b} + i(\mathbf{a} \times \mathbf{b}) \cdot \boldsymbol{\sigma}
\end{gather*}
with vector $\mathbf{a}$ and $\mathbf{b}$.
In components form, the third identity reduce into
\begin{equation*}
    \sigma_i \sigma_j = \delta_{ij} \mathbf{1} + i \epsilon_{ijk} \sigma_k
\end{equation*}
Using the Pauli matrices, we can derive the relation for the spin-$1/2$ operator $\mathbf{S} = \boldsymbol{\sigma}/2$ by multiplying previous relation by $1/4$
\begin{equation*}
    S_iS_j= \frac{1 }{4 }\delta_{ij}\mathbf{1} +\frac{i }{2} \epsilon_{ijk }S_k
\end{equation*}
If $i\neq j$, then 
\begin{equation*}
    S_iS_j = -S_jS_i
\end{equation*}
Since $S_iS_j = i\epsilon_{ijk }S_k/2$ and $S_jS_i= -i \epsilon_{ijk }S_k/2$.
Along with $\mathbf{1}$, the Pauli matrices span the $2 \times  2$ Hermitian matrices.

\subsubsection{Derivation: Pauli matrices identity}
For the first case, we can prove them individually
\begin{align*}
    \sigma_1^2 &= 
    \begin{pmatrix}
        0 & 1 \\
        1 & 0
    \end{pmatrix}
    \begin{pmatrix}
        0 & 1 \\
        1 & 0
    \end{pmatrix}
    =
    \begin{pmatrix}
        1 & 0 \\
        0 & 1
    \end{pmatrix}
    = \mathbf{1}\\
    \sigma_2^2 &= 
    \begin{pmatrix}
        0 & -i \\
        i & 0
    \end{pmatrix}
    \begin{pmatrix}
        0 & -i \\
        i & 0
    \end{pmatrix}
    =
    \begin{pmatrix}
        1 & 0 \\
        0 & 1
    \end{pmatrix}
    = \mathbf{1}\\
    \sigma_3^2 &= 
    \begin{pmatrix}
        1 & 0  \\
        0 & -1
    \end{pmatrix}
    \begin{pmatrix}
        1 & 0  \\
        0 & -1
    \end{pmatrix}
    =
    \begin{pmatrix}
        1 & 0 \\
        0 & 1
    \end{pmatrix}
    = \mathbf{1}
\end{align*}
Next we consider the case 
\begin{equation*}
\sigma_1 \sigma_2 = \begin{pmatrix} 0 & 1 \\ 1 & 0 \end{pmatrix} \begin{pmatrix} 0 & -i \\ i & 0 \end{pmatrix} = \begin{pmatrix} i & 0 \\ 0 & -i \end{pmatrix} = i \sigma_3
\end{equation*}
Instead of considering other case of multiplucation, we can cyclically permutate the index and obtain
\begin{equation*}
\sigma_2 \sigma_3 = i \sigma_1, \quad \sigma_3 \sigma_1 = i \sigma_2
\end{equation*}
or $\sigma_{i}\sigma_j=i\epsilon_{ijk}\sigma_j$.
By combining both case, we have 
\begin{equation*}
    \sigma_i \sigma_j = \delta_{ij} \mathbf{1} + i \epsilon_{ijk} \sigma_k    
\end{equation*}

Using this, we prove the commutator relation
\begin{align*}
    [\sigma_i,\sigma_j] 
    &=   \sigma_i \sigma_j - \sigma_j \sigma_i\\
    &= (\delta_{ij} \mathbf{1} + i \epsilon_{jk} \sigma_k) - (\delta_{ji} \mathbf{1} + i \epsilon_{jik} \sigma_k)\\
    &=i \epsilon_{jk} \sigma_k - i \epsilon_{jik} \sigma_k = i \epsilon_{jk} \sigma_k + i \epsilon_{ijk} \sigma_k \\
    [\sigma_i,\sigma_j] 
    &= 2i \epsilon_{jk} \sigma_k 
\end{align*}
and the anticommutator
\begin{align*}
    \left\{ \sigma_i,\sigma_j \right\}  
    &=   \sigma_i \sigma_j + \sigma_j \sigma_i\\
    &= (\delta_{ij} \mathbf{1} + i \epsilon_{jk} \sigma_k) + (\delta_{ji} \mathbf{1} + i \epsilon_{jik} \sigma_k)\\
    &=\delta_{ij} \mathbf{1} + \delta_{ji} \mathbf{1} = \delta_{ij} \mathbf{1} + \delta_{ij} \mathbf{1} \\
    \left\{ \sigma_i,\sigma_j \right\}  
    &= 2 \delta_{ij} \mathbf{1}
\end{align*}

For the third identity, note that we can write the left term as 
\begin{align*}
    (\mathbf{a} \cdot \boldsymbol{\sigma})(\mathbf{b} \cdot \boldsymbol{\sigma}) 
    &=  a_i b_j \sigma_i \sigma_j =     a_i b_j\delta_{ij} \mathbf{1} + ia_i b_j \epsilon_{ijk} \sigma_k\\
   &=  a_i b_i \mathbf{1} + i (\mathbf{a }\times \mathbf{b })_k \sigma_k \\
    (\mathbf{a} \cdot \boldsymbol{\sigma})(\mathbf{b} \cdot \boldsymbol{\sigma}) 
    &= (\mathbf{a }\cdot \mathbf{b }) + i (\mathbf{a }\times \mathbf{b } )\cdot \boldsymbol{\sigma}
\end{align*}

\subsubsection{Prove of how Pauli matrices span the $2 \times 2$ Hermitian matrices space.}
Start from arbitrary matrices
\begin{equation*}
    H = \begin{pmatrix}
a & b \\
c & d 
\end{pmatrix},
\quad
    H ^\dagger = \begin{pmatrix}
a & c ^* \\
b ^* & d 
\end{pmatrix}, 
\quad a, b, c, d \in \mathbb{C}
\end{equation*}
The Hermitian condition $H = H ^\dagger$ require
\begin{equation*}
    H = \begin{pmatrix}
a & b \\
b^* & d
\end{pmatrix}, \quad a, d \in \mathbb{R}, \quad b \in \mathbb{C}
\end{equation*}
The diagonal must be real since $a = a ^*$ and $d = d ^*$, meanwhile the diagonal can be complex with $c ^* = b$ implies $c  = b ^*$.
We assume that the space can be written as linear combination of 
\begin{equation*}
    H =c_0 \mathbf{1 } + c_1 \sigma_1 + c_2 \sigma_2 + c_3 \sigma_3 =\begin{pmatrix}
a & b \\
b^* & d
\end{pmatrix} 
\end{equation*}
such that 
\begin{equation*}
        \begin{pmatrix}
        c_0&0\\
        0&c_0
    \end{pmatrix}
    + 
    \begin{pmatrix}
        0 & c_1 \\
        c_1 & 0
    \end{pmatrix}
    + 
    \begin{pmatrix}
        0 & -ic_2 \\
        ic_2 & 0
    \end{pmatrix}
    +
    \begin{pmatrix}
        c_3 & 0  \\
        0 & -c_3
    \end{pmatrix}    
    =
    \begin{pmatrix}
c_0 + c_3 & c_1 - ic_2 \\
c_1 + ic_2 & c_0 - c_3
\end{pmatrix}
\end{equation*}
Equating both, we obtain the results for diagonal elements
\begin{equation*}
    \begin{cases}
        c_0 + c_3 = a\\
        c_0 - c_3 = d     
    \end{cases}
    \Rightarrow
    \begin{cases}
        c_0 = \dfrac{a + d}{2}\\
        c_3 = \dfrac{a - d}{2}        
    \end{cases} 
\end{equation*}
and for the off diagonal
\begin{equation*}
    c_1 - i c_2 = b_r + i b_i 
    \quad 
    \Rightarrow 
    \begin{cases}
     c_1 = b_r\\
     c_2 = -b_i        
    \end{cases}
\end{equation*}
All coefficients $c_0, c_1, c_2, c_3$ are uniquely determined, thus $\left\{ \mathbf{1 },\sigma_1 ,\sigma_2,\sigma_3\right\} $ span the space of all $2 \times 2$ Hermitian matrices.

\end{document}

% 
% Unnecessary
% 
% \subsection*{Total Angular Momentum}
% The total angular momentum operator is defined as the sum of the angular momentum and spin $\mathbf{J}=\mathbf{L }+\mathbf{S}$.
% Ladder operator for angular momentum is expressed as
% \begin{equation*}
%     J_{\pm} = J_x \pm i J_y
% \end{equation*}
% Their action on eigenstates $\ket{j,m}$ are
% \begin{equation*}
%     J_{\pm}   |j, m\rangle = \hbar \sqrt{j(j+1) - m(m \pm 1)}   |j, m \pm 1\rangle
% \end{equation*}
% Eigenstate $\ket{j,m}$ is simultaneous eigenstate of the total angular momentum operator $J^2$ with $j $ as the total angular momentum quantum number and $m$ as the magnetic quantum number.
% For given $j$, $m$ can take the value $m=-j,j$ such that there are $2j+1$ degeneracy level.
% The reason it being a degenerate is that the total angular momentum
% \begin{equation*}
%     J^2 |j, m\rangle = \hbar^2 j(j + 1) |j, m\rangle
% \end{equation*}
% are the same, yet the $z$ component
% \begin{equation*}
%     J_z |j, m\rangle = \hbar m |j, m\rangle,
% \end{equation*}
% are different.
% The operator also satisfies
% \begin{equation*}
%     [J_z, J_{\pm}] = \pm \hbar J_{\pm}\qquad[J_+, J_-] = 2\hbar J_z
% \end{equation*}
% 
% \subsubsection*{Derivation.}
% Using the definition of Ladder operator
% \begin{align*}
%     [J_{z}, J_{+}] & =  [J_{z}, J_{x} + i J_{y}] = [J_{z}, J_{x}] + i [J_{z}, J_{y}] = i\hbar J_{y} + i (-i\hbar J_{x}) \\
%                    & =  \hbar (J_{x} + i J_{y})                                                                         \\
%     [J_{z}, J_{+}] & =  \hbar J_{+}                                                                                     \\
%     [J_{z}, J_{-}] & =  [J_{z}, J_{x} - i J_{y}] = [J_{z}, J_{x}] - i [J_{z}, J_{y}] = i\hbar J_{y} - i (-i\hbar J_{x}) \\
%                    & =  -\hbar (J_{x} - i J_{y})                                                                        \\
%     [J_{z}, J_{-}] & =  -\hbar J_{-}
% \end{align*}
% And for other one
% \begin{align*}
%     [J_+, J_-] & =  (J_x + iJ_y)(J_x - iJ_y) - (J_x - iJ_y)(J_x + iJ_y)                     \\
%                & =  J_x^2 - iJ_xJ_y + iJ_yJ_x + J_y^2 - (J_x^2 + iJ_xJ_y - iJ_yJ_x + J_y^2) \\
%                & =  (-iJ_xJ_y + iJ_yJ_x) - (iJ_xJ_y - iJ_yJ_x)                              \\
%                & =  -2iJ_xJ_y + 2iJ_yJ_x                                                    \\
%                & =  2i(J_yJ_x - J_xJ_y)=  2i[J_y, J_x]                                      \\
%     [J_+, J_-] & =  2i(-i\hbar J_z) = 2\hbar J_z
% \end{align*}
% 
% \subsubsection{Derivation: commutator relation (easy mode).}
% First let us derive the basic commutation
% \begin{align*}
%     [\hat{L}_1, \hat{L}_2] & =  [YP_{z} - ZP_{y}, ZP_{x} - XP_{z}]                                        \\
%     [\hat{L}_1, \hat{L}_2] & =  [YP_{z}, ZP_{x}] - [YP_{z}, XP_{z}] - [ZP_{y}, ZP_{x}] + [ZP_{y}, XP_{z}]
% \end{align*}
% Each terms expand as
% \begin{align*}
%     [YP_z,ZP_x]
%      & = Y[P_z,Z]P_x + YZ[P_z,P_x] + [Y,Z]P_zP_x + Z[Y,P_x]P_z \\
%      & = Y[P_z,Z]P_x                                           \\
%      & = -i\hbar Y P_x,                                        \\
%     [YP_z,XP_z]
%      & = Y[P_z,X]P_z + YX[P_z,P_z] + [Y,X]P_zP_z + X[Y,P_z]P_z \\
%      & = 0,                                                    \\
%     [ZP_y,ZP_x]
%      & = Z[P_y,Z]P_x + ZZ[P_y,P_x] + [Z,Z]P_yP_x + Z[Z,P_x]P_y \\
%      & = 0,                                                    \\
%     [ZP_y,XP_z]
%      & = Z[P_y,X]P_z + ZX[P_y,P_z] + [Z,X]P_yP_z + X[Z,P_z]P_y \\
%      & = X[Z,P_z]P_y                                           \\
%      & = i\hbar X P_y
% \end{align*}
% Summing the four results gives
% \begin{align*}
%     [\hat{L}_1, \hat{L}_2] & =  Y[P_{z}, Z]P_{x} + X[Z, P_{z}]P_{y} \\
%                            & =  i\hbar (XP_{y} - YP_{x})            \\
%     [\hat{L}_1, \hat{L}_2] & =  i\hbar \hat{L}_3.
% \end{align*}
% For the posistion relation, we only derive for $X$ case, the other follows immediately
% \begin{flalign*}
%      & [X, \hat{L}_1] = [X, Y P_z - Z P_y] = 0                                      \\
%      & [X, \hat{L}_2] = [X, Z P_x - X P_z] = [X, Z P_x] = Z[X, P_x] = i \hbar Z     \\
%      & [X, \hat{L}_3] = [X, X P_y - Y P_x] = -[X, Y P_x] = -Y[X, P_x] = -i \hbar Y.
% \end{flalign*}
% Same for the momentum
% \begin{flalign*}
%      & [P_{x}, \hat{L}_1] = [P_{x}, Y P_{z} - Z P_{y}] = 0                                                    \\
%      & [P_{x}, \hat{L}_2] = [P_{x}, Z P_{x} - X P_{z}] = -[P_{x}, X P_{z}] = -[P_{x}, X]P_{z} = i \hbar P_{z} \\
%      & [P_{x}, \hat{L}_3] = [P_{x}, X P_{y} - Y P_{x}] = [P_{x}, X P_{y}] = [P_{x}, X]P_{y} = -i \hbar P_{y}  \\
% \end{flalign*}
% 
% Another derivation just for fun
% \begin{align*}
%     [X, L^2] & =  [X, \hat{L}_1^2] + [X, \hat{L}_2^2] + [X, \hat{L}_3^2]                                                    \\
%              & =  0 + \hat{L}_2[X, \hat{L}_2] + [X, \hat{L}_2]\hat{L}_2 + \hat{L}_3[X, \hat{L}_3] + [X, \hat{L}_3]\hat{L}_3 \\
%              & =  i\hbar (\hat{L}_2Z + Z\hat{L}_2 - \hat{L}_3Y - Y\hat{L}_2),
% \end{align*}
% and
% \begin{align*}
%     [P_x, L^2] & =  [P_x, \hat{L}_1^2] + [P_x, \hat{L}_2^2] + [P_x, \hat{L}_3^2]                                                      \\
%                & =  0 + \hat{L}_2[P_x, \hat{L}_2] + [P_x, \hat{L}_2]\hat{L}_2 + \hat{L}_3[P_x, \hat{L}_3] + [P_x, \hat{L}_3]\hat{L}_3 \\
%                & =  i\hbar(\hat{L}_2P_z + P_z\hat{L}_2 - \hat{L}_3P_y - P_y\hat{L}_3).
% \end{align*}
